\documentclass{article}
\usepackage[margin=1in]{geometry}
\usepackage{longtable}
\title{Design Specification\\LA City Sidewalk Assessment Simulation}
\author{Andrew Banagas}
\date{\today}
\begin{document}
\maketitle
\tableofcontents
\newpage

\section{Purpose}
This Design Spec explains every major page, tool, component, package, dependency, and Docker service for the simulated sidewalk assessment system.

\section{Frontend Pages}
\begin{itemize}
  \item Dashboard — Overview of recent scans, processing status
  \item Upload Scan — Interface to import LiDAR/point-cloud files
  \item Scan Review — Viewer for raw scan data and orthoimage outputs
  \item Displacement Results — Display processed displacement CSV and visualizations
  \item Rover Navigation — Manual control and waypoint management UI
  \item Test Logs — View TestRail reports and manual test evidence
  \item Admin Settings — User/role management, model versioning
\end{itemize}

\section{Backend Modules}
\begin{itemize}
  \item File Upload Service — Accept and store raw scan files (LAS/PLY/IMG)
  \item Job Queue Manager — Track processing jobs and status
  \item Results API — Serve processed CSV, images, and metadata
  \item User/Auth Service — Simple role-based access (optional for simulation)
  \item Telemetry Logger — Store rover navigation events and GPS traces
\end{itemize}

\section{ML and Data Processing Modules}
\begin{itemize}
  \item Point Cloud Importer — Read LAS/PLY using Laspy or Open3D
  \item Orthoimage/DEM Generator — Convert point cloud to 2D raster outputs
  \item U-Net Segmentation — TensorFlow/Keras model to identify cracks/joints
  \item Displacement Calculator — Measure vertical and horizontal displacement
  \item CSV Exporter — Format results for BOE review
\end{itemize}

\section{Navigation Modules}
\begin{itemize}
  \item GPS Waypoint Manager — Queue and sequence navigation targets
  \item Obstacle Detection Service — Simulated or real-time detection
  \item Manual Override Handler — User intervention logic
  \item Safe Stop Module — Emergency halt and telemetry snapshot
  \item Telemetry Logger — Record position, battery, events
\end{itemize}

\section{Database and Storage Design}
\begin{itemize}
  \item Metadata DB: PostgreSQL storing scan info, job history, test results
  \item File Storage: Local volume or object store (S3/Azure) for raw/processed files
  \item Model Storage: Versioned U-Net weights at \texttt{/app/storage/models/unet\_model.h5}
\end{itemize}

\section{Dependencies and Libraries}
\subsection{Python 3.10}
\subsection{Backend Stack}
\begin{itemize}
  \item FastAPI or Flask for REST API
  \item python-multipart for file uploads
  \item psycopg2 for PostgreSQL
  \item python-dotenv for environment config
\end{itemize}

\subsection{Data Processing Stack}
\begin{itemize}
  \item numpy — numerical arrays
  \item pandas — data frames and CSV handling
  \item opencv-python-headless — image processing
  \item pillow — image I/O
  \item matplotlib — visualization
  \item scipy — advanced math operations
  \item laspy — read/write LAS files
  \item open3d — point cloud manipulation
\end{itemize}

\subsection{ML Stack}
\begin{itemize}
  \item tensorflow 2.x — deep learning framework
  \item keras — high-level neural network API (bundled with TF)
  \item scikit-learn — preprocessing and metrics
\end{itemize}

\subsection{Frontend}
\begin{itemize}
  \item Node.js 20 LTS
  \item React 18+ — UI framework
  \item Vite — build tool
  \item Tailwind CSS or Material UI — styling
\end{itemize}

\section{Docker Services}
\subsection{frontend}
\begin{itemize}
  \item Base image: \texttt{node:20-alpine}
  \item Runs Vite dev server on port 5173
  \item Mounts source code for live reload
\end{itemize}

\subsection{backend}
\begin{itemize}
  \item Base image: \texttt{python:3.10-slim}
  \item Runs FastAPI/Uvicorn on port 8000
  \item Environment variables: \texttt{DATABASE\_URL}, \texttt{STORAGE\_PATH}
  \item Depends on: db service
\end{itemize}

\subsection{ml-worker}
\begin{itemize}
  \item Base image: \texttt{python:3.10-slim}
  \item Runs background ML processing
  \item Mounts shared storage for models and outputs
  \item Watches job queue and processes scans
\end{itemize}

\subsection{db}
\begin{itemize}
  \item Base image: \texttt{postgres:16}
  \item Database name: sidewalk
  \item User/pass: sidewalk/sidewalk
  \item Persists data in named volume db-data
\end{itemize}

\section{Input and Output File Formats}
\subsection{Inputs}
\begin{itemize}
  \item LAS/PLY: Point cloud format with XYZ + intensity
  \item IMG: Orthoimage (GeoTIFF or similar)
  \item CSV: Waypoint lists (lat, lon, heading, speed)
\end{itemize}

\subsection{Outputs}
\begin{itemize}
  \item Orthoimage: GeoTIFF or PNG at ground resolution
  \item DEM: Elevation raster image
  \item Segmentation Mask: Binary or multi-class PNG
  \item Displacement CSV: Fields = location, vertical\_mm, horizontal\_mm, confidence
  \item Telemetry JSON: GPS trace, events, timestamps
\end{itemize}

\section{Failure Handling}
\begin{itemize}
  \item LiDAR scan import failure: Log error, notify user, mark job as failed
  \item DEM/orthoimage mismatch: Resample or reject misaligned outputs
  \item Missing Docker dependency: Build fails with clear error message
  \item GPS waypoint drift: Log deviation, trigger re-centering or manual override
  \item Obstacle detection false positive: Require manual confirmation before stopping
  \item Safe stop: Halt rover, log telemetry snapshot, await operator reset
\end{itemize}

\section{Security and Access Assumptions}
\begin{itemize}
  \item This is a simulated system for CS 3338 coursework
  \item No authentication required (optional role/login for extensibility)
  \item Field data (images with pedestrians, addresses) not used in simulation
  \item Docker network is local only; no external exposure expected
  \item Future: Implement OAuth2 or JWT for production rover field deployment
\end{itemize}

\end{document}